\section{Slab operators and finite-volume comparison}\label{sec:slab-comparison}

A cubic torus admits the coordinate reflection-positive representations
of the preceding section. A diagonal mirror requires a different period
lattice. To compare the two geometries, we express the functional as a
trace of operators associated with slabs. Reflection gives the adjoint
of a slab operator; repeated reflected copies give its Schatten norm.
The resulting comparison changes the closing vector, and, when needed,
the fermion character. Three such changes connect the cubic torus with
the diagonal torus and with a second coordinate torus of the same volume.
Their pressures differ by a quantity tending to zero. This is the
finite-volume input for diagonal reflection positivity. The slab kernels
and their composition, reflection and trace identities are those of
\cite[Sections~5.1--5.7, Theorem~5.29]{ChamberQED}; the torus comparisons
and pressure argument follow its Section~6.

The same boundary construction will give the coordinate trace used in
\cref{sec:coordinate} to prove Gaussian domination. There the real
external two-cochain enters through unitary electric translations and
magnetic multiplication operators; two choices of time direction
remove its electric and magnetic parts.

All parameters remain fixed. We use the counting Berezin convention
\eqref{eq:reg-berezin}.

\subsection{Local observables and their lifts}\label{subsec:slab-local}

We need local observables that may contain fermion variables and whose
coefficients grow with the curvature. Their definition must make sense
in every flux sector. Let $B=y+[-r,r]^3$, and write $B^{(j)}$ for the
oriented lattice $j$-cells contained in $B$. The fermion sites assigned
to $B$ are
\[
 \mathcal X(B)=\{x\in\mathcal X:
          \pi(x)+[-\tfrac12,\tfrac12]^3\subset B\}.
\]
This support condition includes the photon site at which the fermion
transforms under gauge transformations.

Let $\mathcal G(B)$ consist of the Borel maps
$G:\R^{B^{(1)}}\to\operatorname{Gr}(\mathcal X(B))$ with gauge-invariant
values and with the following growth bound. If
$G(\widehat A)=\sum_\Psi a_\Psi(\widehat A)\Psi$, then, for some
$C<\infty$ and $\varepsilon>0$,
\begin{equation}\label{eq:slab-box-growth}
 |a_\Psi(\widehat A)|\le C
   \exp\left((\tfrac12-\varepsilon)
                    \sum_{\sigma\in B^{(2)}}(d\widehat A)_\sigma^2\right).
\end{equation}
Here gauge invariance means
\begin{equation}\label{eq:slab-box-gauge}
 a_\Psi(\widehat A-d\lambda)
        e^{iq\langle\lambda,\mathfrak q(\Psi)\rangle}
       =a_\Psi(\widehat A)
\end{equation}
for every real function $\lambda$ on the photon sites of $B$;
$\mathfrak q(\Psi)$ counts the $\psi$ generators minus the
$\bar\psi$ generators at fermion sites projecting to each photon site.
We denote the even subspace by $\mathcal G^{\mathrm{ev}}(B)$.

A local lift of a configuration $(A,z)$ on $B$ is a real one-cochain
$\widehat A$ on $B^{(1)}$ satisfying
\begin{equation}\label{eq:slab-local-lift}
 d\widehat A=F\quad\text{on }B^{(2)},\qquad
 e^{iq\widehat A(b)}=e^{iqA(b)}\quad(b\in B^{(1)}).
\end{equation}
We evaluate $G$ at this lift and then impose the fermion identifications
of the torus or cylinder. We now verify that this procedure is Borel
and independent of the lift.

Replace $B$ by its integer box
$\prod_\mu[a_\mu,b_\mu]$, where $a_\mu=\lceil y_\mu-r\rceil$ and
$b_\mu=\lfloor y_\mu+r\rfloor$. If this box has no sites, use the
unique empty lift. Otherwise its axial spanning tree consists of all
links in direction $3$, the links in direction $2$ on $x_3=a_3$, and the
links in direction $1$ on $x_2=a_2$, $x_3=a_3$. Prescribe a one-cochain
on this tree. Its values on the remaining links are determined by
\begin{align}\label{eq:slab-lift-recursion}
 \widehat A_1(x+e_2)&=\widehat A_1(x)+\widehat A_2(x+e_1)
                                  -\widehat A_2(x)-F_{12}(x),
                                  &&x_3=a_3,\notag\\
 \widehat A_1(x+e_3)&=\widehat A_1(x)+\widehat A_3(x+e_1)
                                  -\widehat A_3(x)-F_{13}(x),\\
 \widehat A_2(x+e_3)&=\widehat A_2(x)+\widehat A_3(x+e_2)
                                  -\widehat A_3(x)-F_{23}(x).\notag
\end{align}
First use the first line in increasing $x_2$, then the last two lines
in increasing $x_3$. Each link outside the tree occurs as an unknown
once. The construction is linear in $F$ and the tree values. It gives
the prescribed $13$- and $23$-curvatures and the $12$-curvature on the
bottom face. If $dF=0$, induction over the cubes gives the remaining
$12$-curvatures: on each cube the two closed two-cochains $F$ and
$d\widehat A$ agree on five faces, hence on the sixth.

On the tree choose $\widehat A(b)=\arg_q(e^{iqA(b)})$, where
$\arg_q$ takes values in $[0,r_q)$. This is a Borel choice. Subtract
the recursion for $A$ from \eqref{eq:slab-lift-recursion}. The tree
differences lie in $r_q\Z$, and every subsequent difference is a signed
sum of earlier ones and of a value of $r_qz$. Thus
$\widehat A-A$ is integer-valued in units of $r_q$ on every link. This
proves \eqref{eq:slab-local-lift}, with Borel dependence on $F$ and the
tree phases.

Two local lifts differ by $r_qk$, where $k$ is an integer closed
one-cochain. Sum $k$ along the tree paths from its bottom corner to
define an integer function $\lambda$. Then $d\lambda=k$ on the tree.
The zero-curvature recursion shows that $k-d\lambda=0$ on every link.
Equation \eqref{eq:slab-box-gauge}, applied with $-r_q\lambda$, gives
$a_\Psi(\widehat A+r_qd\lambda)=a_\Psi(\widehat A)$, since
$\langle\lambda,\mathfrak q(\Psi)\rangle$ is integer. The value of $G$
is therefore independent of the lift.

For an open set $U$, let $\mathcal A(U)$ be the linear span of the maps
\begin{equation}\label{eq:slab-admissible-product}
 \widehat A\longmapsto\Phi(d\widehat A)
                   \prod_{j=1}^rG_j(\widehat A|_{B_j^{(1)}}),
\end{equation}
where the boxes $B_j\subset U$ are pairwise disjoint,
$G_j\in\mathcal G^{\mathrm{ev}}(B_j)$, and $\Phi$ is a bounded Borel
function of finitely many curvature values at plaquettes contained in
$U$. A box $B$ embeds in a period group $\Gamma$ when the sets
$B+[-1,1]^3+\gamma$, $\gamma\in\Gamma$, are pairwise disjoint.
Boxes $B_1,\ldots,B_r$ embed disjointly in $\Gamma$ when each
embeds and all sets $B_j+\gamma$, indexed by distinct pairs
$(j,\gamma)$, are pairwise disjoint. A product is read on the
quotient only when its boxes embed disjointly.

It remains to check that equal maps \eqref{eq:slab-admissible-product}
have equal values on configurations. This requires one lift for every
box in a representation. On $\Z^3$, a closed integer two-cochain $z$
has the integer primitive
\begin{align*}
 k_1(x)&=0,\\
 k_2(x)&=\sum_{s=0}^{x_1-1}z_{12}(s,x_2,x_3),\\
 k_3(x)&=\sum_{s=0}^{x_1-1}z_{13}(s,x_2,x_3)
                    +\sum_{s=0}^{x_2-1}z_{23}(0,s,x_3).
\end{align*}
The sums are oriented for negative endpoints. The first two components
of $dk=z$ follow by taking differences. In the third, $dz=0$ changes
the summand into a difference of consecutive $z_{23}$ values, so the
sum telescopes. Therefore $\widehat A=A+r_qk$ is a lift on the whole
lattice. It need not be periodic, even when $(A,z)$ is periodic.

The same conclusion holds on a closed slab. Choose a coordinate
$\mu_0$ absent from the normal $n_\ell$. The slab cell complex is a
planar strip times $\Z$ in direction $\mu_0$.
Let $\mu<\nu$ be the other coordinates, denoted by $(a,r)$.
Write $Q(a,r)$ for the unit square with lower corner $(a,r)$ and
$I_r$ for the interval of $a$ such that $[(a,r),(a,r+1)]$ lies
in the strip. When $I_r$ is nonempty choose $a_r\in I_r$.
On the base sheet set the $\mu$-link values of $k$ to zero and put
\[
 k([(a,r),(a,r+1)]\times\{0\})
       =\sum_{a_r\le s<a}z(Q(s,r)\times\{0\}).
\]
All squares in this oriented sum lie in the strip because $I_r$
is an interval. Set $k=0$ on the $\mu_0$-links. For a horizontal
link $b$ in direction $\nu'\in\{\mu,\nu\}$, extend by
\[
 k(b\times\{j\})=k(b\times\{0\})
       +s_{\nu'}\sum_{0\le j'<j}z(b\times[j',j'+1]),
 \qquad
 s_{\nu'}=\begin{cases}-1,&\nu'<\mu_0,\\+1,&\nu'>\mu_0.\end{cases}
\]
The sums with reversed endpoints have the opposite sign.
The resulting primitive has $dk=z$ on all vertical
plaquettes and on the base sheet. The closed two-cochain $z-dk$ then
has vanishing vertical components; closedness makes its horizontal
component constant between sheets. It vanishes on the base sheet and
hence everywhere. Again the primitive is taken on the lifted strip,
not on its periodic quotient.

Evaluation at this common lift, followed by the quotient homomorphism
on fermion generators, agrees with the separate local evaluations.
It is linear in the map $G$, so it is independent of any representation
of $G$ as a sum of products. Moreover, different box plaquettes have
different representatives on the quotient. Multiplying
\eqref{eq:slab-box-growth} by the Maxwell density therefore leaves
\begin{equation}\label{eq:slab-local-majorant}
 C\exp(-\varepsilon_G\|F\|^2),\qquad
 \varepsilon_G=\min(\tfrac12,\varepsilon_1,\ldots,\varepsilon_r)>0.
\end{equation}
The finite Grassmann expansion of the action has bounded coefficients
for real $A$. The Gaussian sum \eqref{eq:reg-gaussian-sum} now proves
absolute integrability, coefficient by coefficient. This also proves
the required Borel measurability on the quotient.

\subsection{Boundary spaces and slab kernels}\label{subsec:slab-kernels}

Fix a mirror family $\ell$ and a cross-section
$\Xi=(\Gamma_\parallel,\chi_\parallel)$, where
$\Gamma_\parallel\subset\Z^3\cap\{\ell=0\}$ has rank two and systole
at least four. Here and below, the systole of a period lattice
$\Gamma$ is $\min_{\gamma\in\Gamma\setminus\{0\}}|\gamma|_\infty$.
Work on the cylinder $\Z^3/\Gamma_\parallel$. The
admissible mirror positions are
$c\in\frac12|n_\ell|^2\Z$. Write $\mathcal M^j(c)$ for the $j$-cells
contained in $\{\ell=c\}$, $\mathcal X^j(c)$ for those cut through
their relative interiors by this plane, and $\mathcal B^j(c,c')$ for
the $j$-cells contained in $\{c\le\ell\le c'\}$ and in neither
bounding plane. In particular a slab cell may touch either boundary.
All these sets are finite on a bounded slab.

For a chain of admissible mirrors, the cylinder cells partition into
the cells in its mirrors, those cut by them, and those lying between
consecutive mirrors. A cell cut by $c$ has its height range
$[c-|n_\ell|^2/2,c+|n_\ell|^2/2]$. Its non-cut faces lie on one
side of the plane; the other faces are cut by the same plane. This
partition applies also when two mirrors are consecutive. A fermion
bond crosses at most one mirror of such a chain, and no fermion site
lies in a mirror.

A closed-slab configuration $\varpi=(A,z)$ assigns real link values
and integer plaquette values to all cells of the closed slab, with
$dz=0$ on its cubes. We integrate its interior values using
$A(b)\in[0,r_q)$ and counting measure for $z$. Its interior photon
weight is
\begin{equation}\label{eq:slab-interior-weight}
 W_{c,c'}(\varpi)=
    \prod_{\sigma\in\mathcal B^2(c,c')}e^{-F_\sigma^2/2}
    \prod_{Q\in\mathcal B^3(c,c')}\mathbf1[(dz)(Q)=0].
\end{equation}
We use the larger class $\mathcal I_\Xi[c,c']$ of Borel
Grassmann-valued functions of closed-slab configurations, with fermion
support $c<\ell<c'$, invariant under
$(A,z)\mapsto(A+r_qk,z-dk)$, whose coefficients obey
\begin{equation}\label{eq:slab-integrand-growth}
 |a_\Psi(\varpi)|\le C_G
   \exp\left((\tfrac12-\varepsilon_G)
          \sum_{\sigma\in\mathcal B^2(c,c')}F_\sigma^2\right),
 \qquad 0<\varepsilon_G\le\tfrac12.
\end{equation}
The local products supported in the open slab give elements of this
class by the lift construction above. Products in consecutive slabs
remain in the class, with the smaller of the two positive
$\varepsilon$'s. Reflection and translation also preserve the class.

At the mirror $c$, let $\mathcal J(c)$ be the ordered set of crossing
modes from \eqref{eq:reg-crossing-modes}, taken modulo
$\Gamma_\parallel$. Let $\mathcal F(c)$ be the exterior space with
orthonormal basis $e_I$, $I\subset\mathcal J(c)$. Put
$\omega_c e_I=\omega_Ie_I$ and $U_{\mathrm f}e_I=(-1)^{|I|}e_I$.
Translation of the modes includes both the character
$\chi_\parallel$ and the sign required to reorder their product:
\[
 U_ve_I=s_v(I)e_{r_v(I)},\qquad
 \tau_v\zeta_I^{(c)}=s_v(I)\zeta_{r_v(I)}^{(c+\ell(v))}.
\]
The $U_v$ are unitary, satisfy the translation law, and commute with
parity. For an integer translation parallel to the normal, choose the
orders compatibly so that $s_v(I)=1$.

For a Grassmann element $H$ supported in the open slab define
\begin{equation}\label{eq:slab-fermion-matrix}
 N_{c,c'}[H]_{I',I}
    =\int_{(c,c')}\zeta_I^{(c)}H
                        \Theta_{c'}(\zeta_{I'}^{(c')}).
\end{equation}
An operator $\Upsilon$ on $\mathcal F(c)$ is represented at the
mirror by
\[
 \operatorname{Ins}_c[\Upsilon]
       =\sum_{I,J}\Upsilon_{IJ}\Theta_c(\zeta_J^{(c)})\zeta_I^{(c)}.
\]
The particular insertion
$\operatorname{Ins}_c[\omega_c\diag(u_I(A))]$ is the crossing
exponential \eqref{eq:reg-crossing-exponential}. The ordering in these
definitions gives
\begin{equation}\label{eq:slab-matrix-composition}
 N_{c,c''}[H\operatorname{Ins}_{c'}[\Upsilon]H']
             =N_{c',c''}[H']\Upsilon N_{c,c'}[H].
\end{equation}
Indeed, an entry of the left side is a sum of products
$[\zeta_IH\Theta_{c'}\zeta_J]
 [\zeta_KH'\Theta_{c''}\zeta_{I''}]$ in disjoint sets of slab
generators. Berezin factorization gives the entries on the right
without moving any of these factors. Reflecting the integrand reverses
their order and conjugates the Berezin coefficient; with
$v=2(c'-c)n_\ell/|n_\ell|^2$ this gives
\begin{equation}\label{eq:slab-matrix-adjoint}
 N_{c',2c'-c}[\Theta_{c'}H]=U_vN_{c,c'}[H]^*.
\end{equation}

We next include the photon crossing in the boundary data. The cut
incidence maps $d_c$ and $d_c^+$ are the restrictions of the coboundary
to cut faces and to upper faces, as in
\eqref{eq:reg-cut-differentials}. Introduce one angle $\vartheta_Q$
for each cut cube and one integer $k_\sigma$ for each cut plaquette,
and set
\begin{align}\label{eq:slab-dual-measure}
 \varkappa_\sigma&=r_q^{-1}
       ((d_c^*\vartheta)_\sigma-2\pi k_\sigma),\notag\\
 \dd\mu_c^\times&=
      \prod_{Q\in\mathcal X^3(c)}\frac{\dd\vartheta_Q}{2\pi}
      \prod_{\sigma\in\mathcal X^2(c)}
                     \frac{q}{\sqrt{2\pi}}e^{-\varkappa_\sigma^2/2},\\
 e_c((\vartheta,k);\varpi)&=
       e^{i\langle\vartheta,d_c^+z\rangle
                             -i\langle\varkappa,d_c^+A\rangle},\notag\\
 D_c(k,I)&=r_q^{|\mathcal X^1(c)|}
       \mathbf1[d_c^*k+\mathfrak q(I)=0].\notag
\end{align}
Here $\mathfrak q(I)(b)$ is the number of modes of groups $1,2$ over
the cut link $b$, minus the number of groups $3,4$ over it. There are
no cut links at diagonal mirrors. At a coordinate site plane there
are no cut cells at all; its Fourier space is a point.

The one-mirror photon identity proved by the Fourier and Poisson
calculation of the preceding section is
\begin{multline}\label{eq:slab-photon-gluing}
 \sum_{z^\times}\int\dd A^\times\,
       \prod_{\sigma\in\mathcal X^2(c)}e^{-F_\sigma^2/2}
       \prod_{Q\in\mathcal X^3(c)}\mathbf1[(dz)(Q)=0]\,u_I(A^\times)
       \\
 =\int\dd\mu_c^\times\,D_c(k,I)
           e_c((\vartheta,k);\varpi_+)
            \overline{e_c((\vartheta,k);\theta_c^*\varpi_-)}.
\end{multline}
The upper and lower data on the right have the same values on the
mirror. Their upper and reflected lower faces provide the non-cut
faces on the left. Cube Fourier multipliers enforce closedness,
Poisson summation gives the factors $e^{-\varkappa^2/2}$, and compact
cut-link integration gives $D_c(k,I)$. This identity is used in the
direction shown when constructing a kernel and in the opposite
direction when gluing two kernels.

Mirror variables are
$\xi=(\mathfrak m,\vartheta,k,I)$, with
$\mathfrak m=(A_0,z_0)\in[0,r_q)^{\mathcal M^1(c)}
\times\Z^{\mathcal M^2(c)}$. On their space $Y_c$ put
\begin{equation}\label{eq:slab-boundary-measure}
 \dd\mu_c(\xi)=\dd A_0
      \prod_{\sigma\in\mathcal M^2(c)}e^{-F_\sigma(\mathfrak m)^2/2}
                  \dd\mu_c^\times(\vartheta,k)D_c(k,I),
 \qquad \mathcal H(c)=L^2(Y_c,\mu_c).
\end{equation}
The integer and mode variables carry counting measure. This is a
finite positive measure. Summing the mirror plaquette Gaussians gives
$C_q^{|\mathcal M^2(c)|}$, while Poisson summation at fixed angles
gives
\[
 \frac q{\sqrt{2\pi}}\sum_{k_\sigma}e^{-\varkappa_\sigma^2/2}
 \le\sum_{j\in\Z}e^{-r_q^2j^2/2}.
\]
The remaining factors are bounded by
$r_q^{|\mathcal M^1(c)|+|\mathcal X^1(c)|}2^{|\mathcal J(c)|}$.
The induced signed translations $V_v:\mathcal H(c)\to
\mathcal H(c+\ell(v))$ are unitary and satisfy the group law. The
parity operator multiplies the component $I$ by $(-1)^{|I|}$.

Let $S_{(c,c')}(A)$ be the fermion action containing precisely the mass
and bond terms with fermion endpoints in the open slab. For
$G\in\mathcal I_\Xi[c,c']$, define the kernel
\begin{multline}\label{eq:slab-kernel}
 \mathfrak S[G](\xi',\xi)=\sqrt{\omega_{I'}\omega_I}
      \int\dd\varpi^\circ\,W_{c,c'}(\varpi)
      e_c((\vartheta,k);\varpi)
       \overline{e_{c'}((\vartheta',k');\theta_{c'}^*\varpi)}
       \\
       \cdot N_{c,c'}[G(\varpi)e^{-S_{(c,c')}(A)}]_{I',I},
       \qquad \varpi=\mathfrak m\vee\varpi^\circ\vee\mathfrak m'.
\end{multline}
Set the whole integrand to zero where $W_{c,c'}(\varpi)=0$;
$G(\varpi)$ is evaluated only on the closed configurations in its
domain.
The associated integral operator maps $\mathcal H(c)$ to
$\mathcal H(c')$. The square roots of the mode weights belong in the
kernel; the measure \eqref{eq:slab-boundary-measure} contains none of
these weights. This places one full crossing weight at a glued mirror.

\begin{proposition}\label{thm:slab-operators}
Every $\mathfrak S[G]$ is Hilbert--Schmidt and depends linearly on $G$.
If $G$ has constants $C_G,\varepsilon_G$ in
\eqref{eq:slab-integrand-growth}, then
\begin{align}\label{eq:slab-kernel-bound}
 |\mathfrak S[G](\xi',\xi)|
    &\le C_{c,c'}C_GC_q(\varepsilon_G)^{|\mathcal B^2(c,c')|},\notag\\
 \|\mathfrak S[G]\|_2
    &\le\sqrt{\mu_c(Y_c)\mu_{c'}(Y_{c'})}
                C_{c,c'}C_GC_q(\varepsilon_G)^{|\mathcal B^2(c,c')|}.
\end{align}
The constant depends on the fixed cross-section, mirror positions,
$q,M,\rho,w_i$, and chamber base point. For consecutive slabs,
\begin{align}\label{eq:slab-operator-rules}
 \mathfrak S[G']\mathfrak S[G]&=\mathfrak S[GG'],\notag\\
 \mathfrak S[\tau_vG]&=V_v\mathfrak S[G]V_v^{-1},\\
 \mathfrak S[\Theta_{c'}G]&=V_{2(c'-c)n_\ell/|n_\ell|^2}
                                         \mathfrak S[G]^*.\notag
\end{align}
If $\mathbb T=(\Gamma,\chi)$ has cross-section $\Xi$, $g\in\Gamma$
is a closing vector with minimal positive height $p=\ell(g)$, and
$c_0<\cdots<c_r=c_0+p$ are admissible, $r\ge2$, and
$G_j\in\mathcal I_\Xi[c_{j-1},c_j]$ for every $j$, then
\begin{equation}\label{eq:slab-trace}
 Z_{\mathbb T}[G_1\cdots G_r]
    =\operatorname{Tr}_{\mathcal H(c_0)}
       \left(V_g^{-1}U_{\mathrm f}^{(1+\chi(g))/2}
               \mathfrak S[G_r]\cdots\mathfrak S[G_1]\right).
\end{equation}
These identities apply to arbitrary parities. For homogeneous $G$,
$U_{\mathrm f}\mathfrak S[G]=(-1)^{\deg G}\mathfrak S[G]U_{\mathrm f}$.
\end{proposition}

\begin{proof}
The coefficient of the finite Grassmann polynomial
$e^{-S_{(c,c')}(A)}$ is bounded uniformly in real $A$. Each coefficient
in the matrix in \eqref{eq:slab-kernel} is consequently bounded by a
fixed multiple of the majorant for $G$. The phases have modulus one.
After multiplication by the photon weight, the absolute integrand is
bounded by
\[
 C C_G\prod_{\sigma\in\mathcal B^2(c,c')}
                           e^{-\varepsilon_GF_\sigma^2}.
\]
Drop the cube constraints and sum the integers separately. Each sum
is at most $C_q(\varepsilon_G)$; each compact link integral contributes
$r_q$. The mode factors and finite matrix dimension enter $C_{c,c'}$.
This proves the pointwise bound in \eqref{eq:slab-kernel-bound}.
The boundary measures are finite, so the kernel is in
$L^2(Y_{c'}\times Y_c)$, with the stated norm. These bounds also
justify all the gluing integrals below before any cancellation.

For composition, multiply the two kernels and integrate over the common
boundary. Their two square roots give $\omega_{I'}$ for the middle
mode. The two phases and its boundary measure are precisely the right
side of \eqref{eq:slab-photon-gluing}. Applying that identity restores
the photon variables on the cut cells, their Maxwell weights, their
cube constraints, and the phase $u_{I'}$ of the crossing modes. Sum
the middle modes by \eqref{eq:slab-matrix-composition}, with
$\Upsilon=\omega_{c'}\diag(u_{I'})$. This restores the crossing
fermion action. The plaquettes in the middle mirror come from its
boundary measure, and the cell partition restores the full weight on
the combined slab. No cube lies in a mirror, so every closedness
constraint is present once. This is the kernel of $\mathfrak S[GG']$.

Translation acts by a signed relabeling of the modes and an
orientation-preserving relabeling of cells, proving covariance.
For reflection use \eqref{eq:slab-matrix-adjoint}; the two boundary
phases exchange their roles and the interior weight is preserved.
Reflection may reverse a compact integration interval, but integer
shift invariance makes any interval of length $r_q$ equivalent to
$[0,r_q)$. The reflected upper boundary is the translate of the
original lower one by $2(c'-c)n_\ell/|n_\ell|^2$. This gives the last
identity of \eqref{eq:slab-operator-rules}. Parity follows because a
Berezin integral of odd total degree vanishes.

For the trace, composition reduces the claim to two slabs. A product
of two Hilbert--Schmidt kernels is trace class, and its trace is the
absolutely convergent integral of their cyclic product. The closing
operator identifies the upper mirror data with the translate by $g$
of the lower data. Its contribution to this cyclic kernel is
\[
 (-\chi(g))^{|I_0|}s_g(I_0).
\]
Here $s_g$ is the mode relabeling sign; $(-\chi(g))^{|I_0|}$ is the
remaining Grassmann closing sign. It is $1$ for antiperiodic closure
and the fermion parity for periodic closure. To see this sign directly,
translate the lower crossing product to the upper boundary. Each of
its $|I_0|$ generators contributes $\chi(g)$, while moving it across
the adjacent upper product contributes $(-1)^{|I_0|}$ in every term
of even total degree. Their product is $(-\chi(g))^{|I_0|}$.

Apply \eqref{eq:slab-photon-gluing} at the two mirrors. For fixed
photon variables the remaining mode sum is
\[
 \tr_{\mathcal F(c_0)}
    (U_g^{-1}U_{\mathrm f}^{(1+\chi(g))/2}
                       N_2\Upsilon_1N_1\Upsilon_0),
 \qquad \Upsilon_j=\omega_{c_j}\diag(u_I(A)).
\]
Berezin composition and closure replace this trace by the full torus
Berezin integral. The two open-slab actions and the two crossing
actions include each mass and hopping term once. The photon variables
likewise include every torus link and plaquette once, and the cut and
interior cube constraints impose $dz=0$ everywhere. Their measure is
therefore the measure in \eqref{eq:reg-functional}. This proves
\eqref{eq:slab-trace} for even total degree. With odd total degree both
sides vanish, the operator under the trace anticommuting with parity.
Finally, expand nonhomogeneous entries into their even and odd parts.
\end{proof}

\subsection{Reflected folds and Schatten estimates}\label{subsec:slab-folds}

For an operator $Y$ between Hilbert spaces, write
$\|Y\|_s=(\operatorname{Tr}(Y^*Y)^{s/2})^{1/s}$. We use Schatten's
inequality in the following form; see \cite[slides~55--57]{PongeSchatten2025}. If $1\le s_j<\infty$,
$Y_j:H_{j-1}\to H_j$ belongs to the Schatten class of exponent $s_j$,
$H_r=H_0$, and $\sum_js_j^{-1}=1$, then the cyclic product is trace
class and
\begin{equation}\label{eq:slab-schatten-holder}
 |\operatorname{Tr}(Y_r\cdots Y_1)|
                \le\prod_{j=1}^r\|Y_j\|_{s_j}.
\end{equation}
The statement for different spaces follows from the usual inequality
on their orthogonal direct sum: extend each $Y_j$ by zero on all other
summands. The nonzero block of the product is the displayed cyclic
product, and the singular values of each extension are those of $Y_j$.
A unitary closing identification can be absorbed into $Y_r$ without
changing its Schatten norm.

Let $\mathbb T$ have height period $p$ in the family $\ell$. Its
straightened torus $\mathbb T^0$ has the same cross-section and the
closing vector
\begin{equation}\label{eq:slab-straight-vector}
 g^0=\frac p{|n_\ell|^2}n_\ell,
 \qquad \Gamma^0=\Gamma_\parallel\oplus\Z g^0,
 \qquad\chi^0(g^0)=-1.
\end{equation}
We use this torus only when $g^0\in\Z^3$ and its lattice has systole
at least four. Its volume is the volume of $\mathbb T$: both lattices
have the same covolume in the mirror and the same normal height period.

Divide one period into $2k$ slabs of equal height $h'$, with all
$c_j=c_0+jh'$ admissible and $p=2kh'$. For an integrand $G$ in one of
these slabs, form its fold by reflecting it successively in the
bounding planes, taking all $2k$ copies. This is an integrand on
$\mathbb T^0$. Its fermion boxes are disjoint: they lie strictly
between different consecutive mirrors; within each slab the
cross-section embedding separates their translates.

\begin{proposition}\label{prop:slab-fold-holder}
For integrands $G_j\in\mathcal I_\Xi[c_{j-1},c_j]$,
\begin{align}\label{eq:slab-fold-holder}
 Z_{\mathbb T^0}[\operatorname{fold}_\ell G_j]
              &=\|\mathfrak S[G_j]\|_{2k}^{2k},\notag\\
 |Z_{\mathbb T}[G_1\cdots G_{2k}]|
              &\le\prod_{j=1}^{2k}
               Z_{\mathbb T^0}[\operatorname{fold}_\ell G_j]^{1/(2k)}.
\end{align}
In particular every folded factor on the right is nonnegative.
\end{proposition}

\begin{proof}
Fix $j$ and write $a=c_{j-1}$. Put
$Y=\mathfrak S[G_j]:\mathcal H(a)\to\mathcal H(a+h')$, and form
the fold on $[a,a+p]$, with initial boundary $a$ in the trace formula.
Let
$v=2h'n_\ell/|n_\ell|^2$, which is integer because both mirror
positions are admissible. The first reflected pair has product
$V_vY^*Y$ by \eqref{eq:slab-operator-rules}. Its translate by $sv$
has product $V_{(s+1)v}Y^*YV_{sv}^{-1}$. Multiplying all $k$ pairs
telescopes to
\[
 V_{kv}(Y^*Y)^k.
\]
Since $kv=g^0$ and $\chi^0(g^0)=-1$, the closing translation in
\eqref{eq:slab-trace} cancels $V_{g^0}$ and introduces no parity.
The trace is
$\operatorname{Tr}_{\mathcal H(a)}(Y^*Y)^k=\|Y\|_{2k}^{2k}$.
This calculation takes place at the starting boundary of each slab
and proves the first identity for every $G_j$.

For the original torus, the operator
$V_g^{-1}U_{\mathrm f}^{(1+\chi(g))/2}$ is unitary. Apply
\eqref{eq:slab-schatten-holder} to its trace formula, with every
exponent equal to $2k$. The first identity identifies the resulting
norms with folded partition functions. The Hilbert--Schmidt estimate
ensures all operators belong to the required Schatten classes.
\end{proof}

We next apply this inequality to cells. This is the repeated-reflection
chessboard argument of \cite[Section~4]{FILS1978}, with the period
lattice and its fermion character carried through each fold. A frame is a triple of the
mirror families with mutually orthogonal normals $n_1,n_2,n_3$.
The two frames needed below are $(x_1,x_2,x_3)$ and
$(x_1-x_2,x_1+x_2,x_3)$. Choose admissible offsets $o_a$ and an integer
height width $h>0$. The open cells are
\begin{equation}\label{eq:slab-cells}
 Q_\kappa=\{x:o_a+h\kappa_a<\ell^a(x)
                         <o_a+h(\kappa_a+1),\ a=1,2,3\},
 \qquad |Q|=\frac{h^3}{|n_1||n_2||n_3|}.
\end{equation}
The reflections in the planes $\ell^a=o_a+hj$ generate a product of
three infinite dihedral groups. Exactly one element of this group
carries a given open cell to any specified open cell. Its action on
even Grassmann-valued functions is generated by the reflections
$\Theta$ and translations used above. Denote the resulting copy of
$G\in\mathcal A(Q_\kappa)$ in $Q_{\kappa'}$ by $G^{[\kappa']}$.

A rectangular torus in this frame has periods
$\gamma_a=p_an_a/|n_a|^2\in\Z^3$, with $\chi(\gamma_a)=-1$.
If each $p_a\in2h\Z$, the copies fit its period lattice. Let
$N=|\mathbb T|/|Q|$ be the number of cells, and define
\begin{equation}\label{eq:slab-chessboard-norm}
 \operatorname{Fold}_{\mathbb T}G
          =\prod_{\kappa'\bmod\Gamma}G^{[\kappa']},
 \qquad \|G\|_{\mathbb T}
          =Z_{\mathbb T}[\operatorname{Fold}_{\mathbb T}G]^{1/N}.
\end{equation}
This is an unnormalized seminorm. Its defining partition function is
nonnegative, by the first identity of
\eqref{eq:slab-fold-holder} applied to one layer of the folded array.
The value is independent of the chosen representatives of the cells.

\begin{proposition}\label{thm:slab-chessboard}
If a rectangular torus $\mathbb T$ carries the cells
\eqref{eq:slab-cells}, then
\begin{equation}\label{eq:slab-chessboard-rectangular}
 \left|Z_{\mathbb T}\left[\prod_{\kappa}G_\kappa\right]\right|
                       \le\prod_\kappa\|G_\kappa\|_{\mathbb T},
 \qquad G_\kappa\in\mathcal A(Q_\kappa).
\end{equation}
Let $C(a,b,c)$ have coordinate periods $ae_1,be_2,ce_3$,
all with character $-1$.
Let $e_-=e_1-e_2$, $e_+=e_1+e_2$, and let $D(a,b,c)$ have periods
$ae_-,be_+,ce_3$, all with character $-1$. Let $D'(a,b,c)$ have the
same periods and characters, except for $\chi(be_+)=+1$.
For the diagonal frame and $n\in2h\Z$, $n\ge8$, the cells have
representatives
\[
 0\le\kappa_1<n/h,\qquad 0\le\kappa_2<2n/h,
                     \qquad0\le\kappa_3<n/h.
\]
They represent the cells of $C(n,n,n)$, $D'(n/2,n,n)$, and
$D(n/2,n,n)$, and
\begin{equation}\label{eq:slab-chessboard-diagonal}
 \left|Z_{\mathbb T}\left[\prod_\kappa G_\kappa\right]\right|
           \le\prod_\kappa\|G_\kappa\|_{D(n/2,n,n)},
 \quad \mathbb T=C(n,n,n)\ \text{or }D'(n/2,n,n).
\end{equation}
For coordinate cells of width $h$, $a,c\in2h\Z$, $a,c\ge4$, one has
\begin{equation}\label{eq:slab-chessboard-coordinate}
 \left|Z_{D(a,2a,c)}\left[\prod_\kappa G_\kappa\right]\right|
                \le\prod_\kappa\|G_\kappa\|_{C(a,4a,c)}.\qedhere
\end{equation}
\end{proposition}

\begin{proof}
Group the cell integrands into layers perpendicular to the first normal,
and apply \eqref{eq:slab-fold-holder}. Its right side consists of the
arrays obtained by retaining one layer and reflecting it throughout
the straightened torus. Next apply the same inequality to each such
array in direction two, and then in direction three. The reflected
array after these three steps is the full fold of the integrand from
one original cell. The three exponents multiply to
\[
 \frac h{p_1}\frac h{p_2}\frac h{p_3}
       =\frac1N.
\]
Thus the last product consists of one factor
$\|G_\kappa\|_{\mathbb T_3}$ for each original cell, where
$\mathbb T_3$ is the final rectangular torus.

We verify that all three arrays have consistent periods and that the
same cells represent them. If $\Gamma_0$ is the initial lattice, then
straightening successively in the frame gives
\begin{align*}
 \Gamma_1&=(\Gamma_0\cap\{\ell^1=0\})\oplus\Z\gamma_1,\\
 \Gamma_2&=(\Gamma_0\cap\{\ell^1=\ell^2=0\})
                                  \oplus\Z\gamma_1\oplus\Z\gamma_2,\\
 \Gamma_3&=\Z\gamma_1\oplus\Z\gamma_2\oplus\Z\gamma_3,
 \qquad \gamma_a=p_an_a/|n_a|^2.
\end{align*}
The kernel in the second line is a rank-one group along $n_3$; its
positive generator is $\gamma_3$, by the definition of the third
height period. Characters already fixed at $\gamma_1,\gamma_2$ are
retained in subsequent cross-sections, so all three final characters
are $-1$. The brick $0\le\kappa_a<p_a/h$ represents $\Gamma_3$.
It represents $\Gamma_2$ as well; then its section $\kappa_2=0$
represents the second cross-section, and its section $\kappa_1=0$
represents the first one. Restoring the normal coordinate in each
case gives representatives for the preceding lattice. These groups
all have the same covolume.

We also check the stronger box embedding condition used in the local
classes. If $B=y+[-r,r]^3$ lies in an open cell, then
$2|n_a|^2r<h$ for every $a$, by taking the range of $\ell^a$ on $B$.
For a nonzero period $\gamma$, at least one height is nonzero and
$|\ell^a(\gamma)|\ge2h$. Consequently
$|\gamma|_\infty>4r$, and the systole also gives
$|\gamma|_\infty\ge4$. These inequalities imply
$|\gamma|_\infty>2r+2$: use the second if $r<1$, the first if
$r\ge1$. Thus translates of $B+[-1,1]^3$ are disjoint. Boxes in
different cell orbits cannot meet after period translations, since
the open cells themselves are disjoint. It follows that every layer
belongs to its slab class. This justifies the three applications.

For a rectangular torus, every straightening leaves the torus unchanged,
which proves \eqref{eq:slab-chessboard-rectangular}. For a cubic torus
with even side $n$, the diagonal operations are
\begin{equation}\label{eq:slab-three-untwistings}
 C(n,n,n)\xrightarrow{\ell_-}D'(n/2,n,n)
          \xrightarrow{\ell_+}D(n/2,n,n),
 \qquad D(n/2,n,n)\xrightarrow{x_1}C(n/2,2n,n).
\end{equation}
For the first operation the cross-section is generated by $ne_+,ne_3$,
with characters $+1,-1$; its height period is $n$, and straightening
replaces the closing vector $ne_1$ by $(n/2)e_-$. For the second, the
cross-section is generated by $(n/2)e_-,ne_3$, both with character
$-1$. Its closing vector $ne_+$ already points along the normal,
and the operation changes its character from $+1$ to $-1$. For the
third, the cross-section is generated by $2ne_2,ne_3$, both with
character $-1$; the height period is $n/2$ and the straight closing
vector is $(n/2)e_1$. The lattices have systole at least four.

The first two diagonal operations, followed by the unchanged $x_3$
operation, have height periods $(n,2n,n)$. Hence they give
\eqref{eq:slab-chessboard-diagonal}; starting at $D'$ instead, use
the order $(\ell_+,x_3,\ell_-)$ to obtain the same final fold.
Likewise, starting at $D(a,2a,c)$ and using coordinate directions
gives periods $(a,4a,c)$ and ends at $C(a,4a,c)$, proving
\eqref{eq:slab-chessboard-coordinate}.
\end{proof}

\subsection{Pressure and the comparison error}\label{subsec:slab-pressure}

Write $C(n)=C(n,n,n)$ and
\[
 f(n_1,n_2,n_3)=\frac1{n_1n_2n_3}\log Z_{C(n_1,n_2,n_3)}[1].
\]
The mass and charge in this function are held fixed. The following
error will compare local expectations in the cubic and diagonal
geometries.

\begin{proposition}\label{prop:slab-pressure}
The pressure $f$ is nonincreasing in each side, and
\[
 f(n_1,n_2,n_3)\longrightarrow f_\infty:=\inf f\in\R
                     \quad\text{as }\min_\mu n_\mu\longrightarrow\infty.
\]
For every even $n\ge8$,
\begin{equation}\label{eq:slab-partition-chain}
 Z_{C(n)}[1]\le Z_{D'(n/2,n,n)}[1]
       \le Z_{D(n/2,n,n)}[1]\le Z_{C(n/2,2n,n)}[1].
\end{equation}
With $\Delta(n)=f(n/2,2n,n)-f(n,n,n)$,
\begin{equation}\label{eq:slab-pressure-error}
 0\le\Delta(n)\le f(n/2,n/2,n/2)-f_\infty\longrightarrow0.
\end{equation}
The logarithms of any two partition functions in
\eqref{eq:slab-partition-chain} differ in absolute value by at most
$n^3\Delta(n)$.
\end{proposition}

\begin{proof}
Fix the cross-section with sides $n_2,n_3$ and let
$Y=\mathfrak S[1]:\mathcal H(0)\to\mathcal H(1/2)$ be the coordinate
half-step. The first identity of \eqref{eq:slab-fold-holder}, applied
to $2n_1$ half-slabs, gives
\begin{equation}\label{eq:slab-pressure-trace}
 Z_{C(n_1,n_2,n_3)}[1]=\operatorname{Tr}(Y^*Y)^{n_1}.
\end{equation}
The positive trace-class operator $Y^*Y$ is independent of $n_1$ and
is nonzero, since the partition function is positive. Consequently
\[
 f(n_1,n_2,n_3)=\frac1{n_2n_3}\log\|Y^*Y\|_{n_1}.
\]
Schatten norms decrease as their exponent increases, so $f$ decreases
in $n_1$. Permuting the time direction proves the other two claims.
The same trace identity shows separate convexity of $\log Z$ in the
side lengths: H\"older's inequality gives
$\operatorname{Tr}A^{(1-t)s+ts'}\le
 (\operatorname{Tr}A^s)^{1-t}(\operatorname{Tr}A^{s'})^t$ for $A=Y^*Y$.

The lower partition bound \eqref{eq:reg-partition-bounds} gives the
volume-independent floor
\[
 f\ge\log C_*+3\log r_q-6r_q^2.
\]
Together with the upper bound in the same equation, it makes
$f_\infty$ finite. Monotonicity now gives convergence in any manner
with all sides tending to infinity: choose a finite triple at which
$f<f_\infty+\delta$ and then exceed its three sides.

Apply \eqref{eq:slab-fold-holder} with constant integrands to each of
the three operations \eqref{eq:slab-three-untwistings}. A fold of
the constant is the constant, so $Z_{\mathbb T}[1]\le Z_{\mathbb T^0}[1]$.
This proves the chain. Each torus in it has volume $n^3$, so its first
and last logarithms differ by $n^3\Delta(n)$. Side monotonicity gives
$f(n/2,2n,n)\le f(n/2,n/2,n/2)$, while
$f(n,n,n)\ge f_\infty$, proving \eqref{eq:slab-pressure-error}.
All intermediate logarithms lie between the endpoints, which proves
the last assertion.
\end{proof}

Thus the estimate concerns the logarithm per unit volume. It does not
give a bounded ratio of partition functions: the retained bound is
$e^{n^3\Delta(n)}$, with $\Delta(n)\to0$. The local comparison
divides the logarithmic error bound $n^3\Delta(n)$ by the number of
copies of the observable and by the perturbation strength.

